%-------------------------------------------------------------------------
%	QUOTATION PAGE
%--------------------------------------------------------------------------

\quotepage
{The Universe is under no obligation to make sense to you.}
{\textbf{Neil deGrasse Tyson}, \textit{Astrophysics for People in a Hurry (2017).}}
{Nothing in life is to be feared; it is only to be understood.}
{\textbf{Marie Curie}, \textit{Madame Curie: A Biography (1937).}}



%-------------------------------------------------------------------------
%	ACKNOWLEDGEMENTS PAGE
%-------------------------------------------------------------------------
\begin{acknowledgements}

I'm deeply grateful to my husband for his unwavering support, patience, and encouragement throughout this journey. His love and understanding have been a constant source of strength, and I am truly thankful for his presence in my life.

Also, I thank my father-in-law and mother-in-law for their emotional presence, vibrating, fearing, and feeling each step of this journey with me while continued to encourage and support me.

I need to express my sincere gratitude to my mother, who, even though not physically present, has been a guiding light in my life. Her love, wisdom, and values continue to inspire me every day, and I carry her memory with me as I face the challenges and triumphs of this journey. I hope to honor her legacy by doing good to all who cross my path. Without her love and presence, I would not have made it this far.

I am grateful to my father, my sister, and my brother for being who they are — incredible people I am privileged to have in my life — and for their unconditional love. Their existence is fundamental to me staying strong, moving forward, and pursuing my dreams.

I would like to express my sincere love to my cats Haroldo and Luna, that were by my side every day during this work. Filling my life with joy, love, and companionship, they have been a source of comfort and emotional support, making the challenges of this journey more bearable and the achievements even sweeter.

I am also grateful to my friends and colleagues for their clarifications, companionship, encouragement, and emotional support.

Finally, I would like to express my deepest gratitude to Eduardo de Matos Rodrigues, Hélder Filipe Pinto de Oliveira e Tânia Pereira for their invaluable guidance, support, and encouragement throughout this research journey. Their expertise and insights have been instrumental in shaping this thesis. 


\end{acknowledgements}

\addvspacetoc{0.3cm} % Add a gap in the Contents, for aesthetics


%-------------------------------------------------------------------------
%	ABSTRACT PAGE
%-------------------------------------------------------------------------
\begin{abstract}

Metal artifacts remain a common and clinically relevant cause of image degradation in\acrfull{ct}, estimated to affect about 21\% of examinations and potentially obscuring or mimicking pathology. This dissertation develops and rigorously evaluates a reproducible deep-learning \textit{pipeline} for image-domain metal artifact reduction, comparing a conventional U-Net with a dual-supervision variant that adds an auxiliary artifact-mask prediction output. The\acrfull{ctmar} benchmark dataset from the\acrfull{aapm} was used for internal development, and samples from datasets in\acrfull{tcia} were used for external evaluation. The pre-specified primary hypothesis---that reconstruction improves structural similarity specifically in the region directly affected by the artifact while keeping clean tissue comparatively unchanged---was confirmed using an\acrfull{ssim}-based metric applied separately to the artifact\acrfull{roi} and the unaffected region. A direct paired comparison showed a small but significant global advantage for the single-supervision model and a small, also significant Clean-Zone preservation advantage for the dual-supervision model, indicating a correction--preservation trade-off rather than a clear superiority of either architecture. Based on the criterion of greatest improvement in the region affected by the artifact—while maintaining minimal proportional change in the unaffected region—the single-supervision model was selected for external evaluation. An evaluation by five physicians found that the model's output better preserved the anatomical region underlying the artifact and reduced the perceived intensity of the artifact; the physicians also preferred the model's output for clinical interpretation. Given the study design, these results are interpreted cautiously as supporting—though not conclusive—evidence of anatomical integrity and perceived clinical benefit. \\
\\
\noindent\textbf{Keywords:} \keywordnames
\end{abstract}

%-------------------------------------------------------------------------
%	ABSTRACT PAGE (PORTUGUESE)
%-------------------------------------------------------------------------
\begin{Resumo}
\begin{otherlanguage}{portuguese}
Artefatos metálicos permanecem uma causa comum e clinicamente relevante de degradação de imagem em\acrfull{ct}, estimados em cerca de 21\% dos exames, podendo esconder ou simular patologias. Esta dissertação desenvolve e avalia rigorosamente um \textit{pipeline} reprodutível de aprendizagem profunda para redução de artefatos metálicos no domínio da imagem, comparando uma U-Net convencional com uma variante de dupla supervisão que acrescenta uma saída auxiliar de predição da máscara de artefato. Utilizando o \textit{dataset} de referência\acrfull{ctmar} da\acrfull{aapm} e amostras de Datasets do\acrfull{tcia} para avaliação externa. A hipótese primária, definida à partida, de que a reconstrução melhora a similaridade estrutural especificamente na região diretamente afetada pelo artefato, mantendo o tecido limpo comparativamente inalterado, foi confirmada usando uma métrica baseada em\acrfull{ssim} aplicada na\acrfull{roi} do artefato e na região não afetada separadamente. Uma comparação pareada direta mostrou uma pequena mas significativa vantagem global do modelo de supervisão única e uma pequena e também significativa vantagem de preservação na região limpa do modelo de dupla supervisão, indicando um compromisso entre correção e preservação, e não uma superioridade clara de nenhuma das arquiteturas. Com base no critério de maior melhoria da região afetada pelo arfato enquanto se mantém pouca alteração proporcional na região não afetada, o modelo de supervisão única foi selecionado para seguir para uma avaliação externa. A avaliação feita por cinco médicos, considerou que a saída do modelo preservava melhor a região anatómica subjacente ao artefato, reduzia a intensidade percebida do artefato e preferiram a saída do modelo para interpretação clínica. Estes resultados são interpretados de forma cuidadosa, como evidência de suporte, mas não conclusiva, da segurança anatómica externa e do benefício clínico percebido devido ao desenho do estudo do presente trabalho. \\
\\
\noindent\textbf{Palavras-chave:} \keywordnamesPT
\end{otherlanguage}
\end{Resumo}

%-------------------------------------------------------------------------
%	LIST OF CONTENTS/FIGURES/TABLES
%-------------------------------------------------------------------------

\addvspacetoc{0.3cm}

\tableofcontents % Write out the Table of Contents

\listoffigures % Write out the List of Figures

\listofgraphics % Write out the List of Graphics

\listoftables % Write out the List of Tables

%\addvspacetoc{0.3cm}

%-------------------------------------------------------------------------
%	PHYSICAL CONSTANTS/OTHER DEFINITIONS
%-------------------------------------------------------------------------

%\begin{listofcontants}
%	\const{My little ponny test of magical rainbow}{$mn/mp$}
%    {$2.997\ 924\ 58\times10^{8}\ \mbox{ms}^{-\mbox{s}}$}
%   \const{Vaccuum permeability test of magical rainbow for a specific case of
%   condensed matter physics}
%   {$\epsilon_0$}{$2.997\ 924\ 58\times10^{8}\ \mbox{ms}^{-\mbox{s}}$}
%	\const{Speed of Light test of magical rainbow}{$c$}
%    {$2.997\ 924\ 58\times10^{8}\ \mbox{ms}^{-\mbox{s}}$}
%\end{listofcontants}


%-------------------------------------------------------------------------
%	SYMBOLS
%-------------------------------------------------------------------------

%\begin{listofsymbols}
%	\symb{$F_{\mu\nu}$}{Maxwell tensor}{F}
%	\symb{$a$}{distance}{m}
%	\\
%	\symb{$\omega$}{angular frequency}{rads$^{-1}$}
%\end{listofsymbols}


%-------------------------------------------------------------------------
%	NOTATION
%-------------------------------------------------------------------------

% \newcommand\notationname{Notation and Conventions}
% \addtotoc{\notationname}
% \fancyhead[LO]{\textsc{\notationname}}

% \input{Notation}



%-------------------------------------------------------------------------
%	ABBREVIATIONS
%-------------------------------------------------------------------------

% Adicionar todas as abreviações definidas (mesmo as não usadas no texto)
\glsaddall

\printglossary[type=\acronymtype,title=Acronyms]


%-------------------------------------------------------------------------
%	DEDICATORY
%-------------------------------------------------------------------------

\begin{dedicatory}
	This thesis is dedicated to my beloved late mother, who would have been filled with pride and genuine happiness with each step of this achievement. 
\end{dedicatory}
